% Môn: Vật lý
% Lớp: 10
% Chương: Chương I. Mở đầu
% Bài: L10C1 Bài 2. Các quy tắc an toàn trong phòng thực hành Vật lí · Vận dụng quy tắc an toàn khi sử dụng các thiết bị đo điện, nhiệt và thủy tinh.
% Số câu: 185
% App đọc file này trực tiếp — sửa rồi Commit trên GitHub, bấm Đồng bộ GitHub .tex

% L10C1 Bài 2. Các quy tắc an toàn trong phòng thực hành Vật lí



% ===== Câu 27 =====
% ID: L10C1B2-108-TN
% Mức: NB
\dangbt{Vận dụng quy tắc an toàn khi sử dụng các thiết bị đo điện, nhiệt và thủy tinh.}
\begin{ex}
% ID: L10C1B2-108-TN
% Mức: NB
Thao tác nào sau đây là đúng quy chuẩn và đảm bảo an toàn khi tháo phích cắm điện của thiết bị thí nghiệm ra khỏi ổ cắm?
\choice
{Dùng một tay cầm chặt phần dây dẫn điện và giật mạnh ra khỏi ổ cắm}
{Dùng hai tay kéo mạnh cả dây dẫn và phích cắm cùng một lúc}
{\True Cầm vào phần phích cắm bằng nhựa khô và rút dứt khoát ra khỏi ổ cắm}
{Dùng một vật bằng kim loại lách vào khe ổ cắm để bẩy phích cắm ra}
\loigiai{
Tuyệt đối không giật rút trực tiếp dây dẫn điện vì có thể làm đứt lõi dây hoặc hỏng ổ cắm gây rò rỉ điện. Thao tác đúng là cầm vào phần vỏ nhựa khô của phích cắm để rút thẳng dứt khoát.
}
\end{ex}

% ===== Câu 28 =====
% ID: L10C1B2-109-TN
% Mức: NB
\begin{ex}
% ID: L10C1B2-109-TN
% Mức: NB
Khi lắp mạch điện thí nghiệm, quy tắc an toàn quan trọng nhất trước khi nối dây là gì?
\choice
{Để nguồn điện ở điện áp lớn nhất}
{\True Ngắt nguồn điện trước khi lắp hoặc thay đổi mạch}
{Chạm tay kiểm tra hai cực của nguồn}
{Dùng dây dẫn bất kì, kể cả dây bị tróc vỏ}
\loigiai{
Trước khi lắp hoặc thay đổi mạch điện phải ngắt nguồn để tránh điện giật, chập điện và hư hỏng thiết bị.
}
\end{ex}

% ===== Câu 31 =====
% ID: L10C1B2-110-TN
% Mức: NB
\begin{ex}
% ID: L10C1B2-110-TN
% Mức: NB
Để vệ sinh và loại bỏ bụi bẩn bám trên bề mặt các linh kiện quang học nhạy cảm như thấu kính hội tụ, học sinh nên thực hiện thao tác nào dưới đây?
\choice
{Dùng vải thô ráp thấm cồn chà xát thật mạnh lên bề mặt kính}
{\True Dùng bóng cao su thổi sạch bụi và lau nhẹ bằng khăn mềm khô chuyên dụng}
{Rửa trực tiếp thấu kính dưới vòi nước máy rồi phơi dưới ánh nắng gay gắt}
{Dùng giấy nhám mịn mài nhẹ bề mặt thấu kính để loại bỏ vết mốc bám}
\loigiai{
Việc vệ sinh thấu kính phải tuân thủ nguyên tắc nhẹ nhàng: dùng bóng xịt khí thổi sạch hạt bụi lớn trước để tránh làm xước kính, sau đó dùng khăn mềm chuyên dụng lau nhẹ nhàng để không làm trầy xước lớp màng quang học.
}
\end{ex}

% ===== Câu 32 =====
% ID: L10C1B2-112-TN
% Mức: NB
\begin{ex}
% ID: L10C1B2-112-TN
% Mức: NB
Khi nhìn thấy biển cảnh báo bề mặt nóng gần một thiết bị thí nghiệm, học sinh cần làm gì?
\choice
{Chạm tay kiểm tra nhiệt độ của thiết bị}
{\True Không chạm trực tiếp và dùng dụng cụ bảo hộ khi cần thao tác}
{Đổ nước lạnh ngay lên thiết bị}
{Di chuyển thiết bị bằng tay trần}
\loigiai{
Biển cảnh báo bề mặt nóng cho biết thiết bị có thể gây bỏng. Cần tránh chạm trực tiếp và dùng dụng cụ bảo hộ phù hợp.
}
\end{ex}

% ===== Câu 33 =====
% ID: L10C1B2-113-TN
% Mức: NB
\begin{ex}
% ID: L10C1B2-113-TN
% Mức: NB
Hành vi nào sau đây cực kì nguy hiểm và bị nghiêm cấm khi học sinh sử dụng nguồn phát tia laser trong thí nghiệm quang học?
\choice
{Chiếu tia laser song song với mặt bàn thí nghiệm phẳng}
{\True Nhìn trực tiếp bằng mắt thường vào đường truyền của luồng tia laser phát ra từ nguồn}
{Chiếu tia laser đi qua thấu kính phân kì để quan sát chùm tia ló}
{Tắt nguồn phát tia laser ngay sau khi kết thúc các thao tác đo đạc}
\loigiai{
Tia laser có mật độ năng lượng cực cao, nếu chiếu trực tiếp vào mắt sẽ lập tức gây bỏng võng mạc, tổn thương các tế bào nhạy sáng và có nguy cơ cao dẫn đến mù lòa vĩnh viễn.
}
\end{ex}

% ===== Câu 34 =====
% ID: L10C1B2-114-TN
% Mức: TH
\begin{ex}
% ID: L10C1B2-114-TN
% Mức: TH
Trước khi cắm phích điện của máy biến áp vào nguồn điện lưới trong phòng thực hành, người vận hành bắt buộc phải thực hiện thao tác nào để đảm bảo an toàn cho các thiết bị điện trong mạch tải?
\choice
{Bật công tắc nguồn của máy biến áp lên mức lớn nhất}
{\True Vặn nút điều chỉnh điện áp đầu ra về mức thấp nhất}
{Kết nối trực tiếp các thiết bị tải vào ngõ ra của máy biến áp}
{Điều chỉnh đột ngột thang đo điện áp đầu ra khi đang có tải}
\loigiai{
Khi sử dụng máy biến áp, người vận hành bắt buộc phải vặn nút điều chỉnh điện áp đầu ra về mức thấp nhất trước khi bật nguồn điện, sau đó mới tăng dần điện áp lên trị số cần thiết để tránh hiện tượng sốc điện đột ngột gây hỏng hóc các thiết bị trong mạch tải.
}
\end{ex}

% ===== Câu 35 =====
% ID: L10C1B2-115-TN
% Mức: TH
\begin{ex}
% ID: L10C1B2-115-TN
% Mức: TH
Khi dùng kẹp gỗ để giữ ống nghiệm đang đun, mục đích chính là gì?
\choice
{Làm tăng tốc độ bay hơi của chất lỏng}
{\True Tránh bỏng tay và giữ ống nghiệm chắc chắn}
{Làm giảm khối lượng ống nghiệm}
{Làm cho số đo nhiệt độ chính xác hơn}
\loigiai{
Kẹp gỗ giúp giữ ống nghiệm ở khoảng cách an toàn, tránh tiếp xúc trực tiếp với vật nóng.
}
\end{ex}

% ===== Câu 36 =====
% ID: L10C1B2-116-TN
% Mức: TH
\begin{ex}
% ID: L10C1B2-116-TN
% Mức: TH
Để đo cường độ dòng điện một chiều chưa rõ giá trị chạy qua một đoạn mạch sử dụng nguồn pin, học sinh sử dụng một ampe kế có hai thang đo là $0,6\,\mathrm{A}$ và $3\,\mathrm{A}$. Hãy đề xuất cách chọn thang đo an toàn cho ampe kế trước khi đóng công tắc mạch điện.
\choice
{Chọn thang đo nhỏ nhất $0,6\,\mathrm{A}$ để có độ chính xác cao nhất}
{Chọn thang đo bất kì mà không cần quan tâm đến trị số}
{\True Chọn thang đo lớn nhất $3\,\mathrm{A}$ rồi hạ dần xuống nếu cần}
{Chọn thang đo lớn nhất và mắc song song ampe kế với bóng đèn}
\loigiai{
Để tránh dòng điện thực tế trong mạch vượt quá giới hạn chịu đựng của thang đo làm hỏng cuộn dây đo, quy tắc an toàn bắt buộc phải chọn thang đo lớn nhất là $3\,\mathrm{A}$ trước tiên, sau đó mới hạ dần xuống thang đo nhỏ hơn nếu dòng điện thực tế nằm trong giới hạn của thang đo đó.
}
\end{ex}

% ===== Câu 37 =====
% ID: L10C1B2-118-TN
% Mức: TH
\begin{ex}
% ID: L10C1B2-118-TN
% Mức: TH
Trong quá trình thực hiện thí nghiệm về mạch điện, một học sinh phát hiện dây dẫn nối từ nguồn điện đến thiết bị bị hở vỏ cách điện và chạm vào nhau, gây ra tia lửa điện nhỏ và mùi khét. Hành động an toàn và cấp thiết nhất mà học sinh cần làm là gì?
\choice
{\True Ngắt ngay nguồn điện chính của toàn bộ mạch thí nghiệm và thông báo cho giáo viên.}
{Dùng băng dính cách điện quấn lại chỗ dây bị hở để tiếp tục thí nghiệm.}
{Tiếp tục quan sát xem hiện tượng có tự hết không trước khi báo cáo.}
{Rút phích cắm của thiết bị đang sử dụng ra khỏi ổ cắm.}
\loigiai{
Khi phát hiện sự cố chập điện, tia lửa điện và mùi khét, đây là dấu hiệu của nguy hiểm cháy nổ hoặc điện giật. Hành động an toàn và cấp thiết nhất là phải ngắt ngay nguồn điện chính để loại bỏ nguy cơ, sau đó thông báo cho giáo viên để được xử lý đúng cách. Việc tự ý sửa chữa, tiếp tục quan sát hoặc chỉ rút phích cắm của thiết bị đều không đảm bảo an toàn tuyệt đối và có thể gây hậu quả nghiêm trọng hơn.
}
\end{ex}

% ===== Câu 38 =====
% ID: L10C1B2-119-TN
% Mức: TH
\begin{ex}
% ID: L10C1B2-119-TN
% Mức: TH
Trong một buổi thực hành Vật lí, bạn An đang sử dụng một nguồn điện một chiều có hiệu điện thế 24 $V$ để cấp cho mạch điện. Bạn Bình, do tò mò, định chạm tay vào các đầu nối dây điện đang có điện. Hành động nào sau đây của An là đúng đắn và hiệu quả nhất để ngăn chặn nguy hiểm cho Bình?
\choice
{La lớn để cảnh báo Bình dừng lại.}
{Chờ giáo viên đến để xử lý tình huống.}
{\True Nhanh chóng ngắt nguồn điện chính của mạch thí nghiệm.}
{Dùng một vật cách điện (như thước nhựa) để đẩy tay Bình ra.}
\loigiai{
Để ngăn chặn nguy hiểm điện giật cho Bình, hành động đúng đắn và hiệu quả nhất là loại bỏ nguồn điện. Việc nhanh chóng ngắt nguồn điện chính sẽ ngay lập tức làm mất điện trong mạch, loại bỏ hoàn toàn nguy cơ điện giật.
}
\end{ex}

% ===== Câu 42 =====
% ID: L10C1B2-120-TN
% Mức: TH
\begin{ex}
% ID: L10C1B2-120-TN
% Mức: TH
Bạn Minh đang thực hiện thí nghiệm đun nóng chất lỏng bằng bếp điện. Trong lúc thí nghiệm, Minh phát hiện dây nguồn của bếp điện bị sờn, lộ rõ lõi đồng bên trong. Hành động đúng đắn và an toàn nhất mà Minh cần làm là gì?
\choice
{Tiếp tục sử dụng bếp vì nó vẫn đang hoạt động bình thường.}
{Dùng băng dính cách điện quấn lại chỗ dây bị sờn và tiếp tục thí nghiệm.}
{Rút phích cắm, sau đó tự mình cố gắng sửa chữa dây nguồn bị hỏng.}
{\True Ngừng sử dụng bếp điện ngay lập tức, rút phích cắm ra khỏi ổ điện và báo cáo tình trạng này cho giáo viên.}
\loigiai{
ây điện bị sờn, lộ lõi đồng là một nguy cơ nghiêm trọng về điện giật và chập cháy. Hành động an toàn nhất là phải ngừng sử dụng thiết bị ngay lập tức để loại bỏ nguy hiểm, rút phích cắm để đảm bảo không còn dòng điện chạy qua, và báo cáo cho giáo viên để thiết bị được kiểm tra và sửa chữa bởi người có chuyên môn. Việc tiếp tục sử dụng, tự ý sửa chữa tạm thời hoặc tự sửa chữa đều tiềm ẩn rủi ro rất cao và không được phép trong phòng thí nghiệm.
}
\end{ex}

% ===== Câu 43 =====
% ID: L10C1B2-121-TN
% Mức: TH
\begin{ex}
% ID: L10C1B2-121-TN
% Mức: TH
Khi thực hiện thí nghiệm với các dung dịch hóa chất, một học sinh vô tình làm đổ một ít dung dịch axit loãng lên tay. Theo quy tắc an toàn trong phòng thực hành, hành động nào sau đây là ĐÚNG đắn nhất để xử lý tình huống này?
\choice
{Lau khô ngay bằng giấy ăn và tiếp tục làm thí nghiệm.}
{\True Rửa ngay vùng bị dính axit dưới vòi nước sạch trong vài phút và báo cáo giáo viên.}
{Bôi một ít dung dịch kiềm mạnh lên vùng bị dính để trung hòa.}
{Chờ giáo viên đến và tự xử lý khi có thời gian.}
\loigiai{
Khi tiếp xúc với axit, việc đầu tiên và quan trọng nhất là phải rửa ngay vùng bị dính dưới vòi nước sạch để loại bỏ bớt axit và giảm thiểu tác hại. Sau đó, cần báo cáo cho giáo viên để có biện pháp xử lý tiếp theo. Lau khô bằng giấy ăn không có tác dụng trung hòa hay làm sạch. Bôi dung dịch kiềm mạnh có thể gây bỏng nặng hơn do phản ứng trung hòa tỏa nhiệt. Chờ đợi là không nên vì axit có thể gây tổn thương da.
}
\end{ex}

% ===== Câu 44 =====
% ID: L10C1B2-122-TN
% Mức: TH
\begin{ex}
% ID: L10C1B2-122-TN
% Mức: TH
Khi sử dụng các dụng cụ đo điện như ampe kế, vôn kế trong phòng thực hành, quy tắc an toàn nào sau đây là quan trọng nhất để tránh làm hỏng thiết bị và đảm bảo an toàn cho người sử dụng?
\choice
{Luôn kết nối ampe kế song song với mạch điện cần đo cường độ dòng điện.}
{Luôn kết nối vôn kế nối tiếp với mạch điện để đo hiệu điện thế.}
{\True Đảm bảo thang đo của thiết bị phù hợp với giới hạn đo dự kiến của đại lượng cần đo.}
{Chỉ sử dụng các dây dẫn có tiết diện nhỏ để tiết kiệm vật liệu.}
\loigiai{
Việc chọn đúng thang đo cho ampe kế hoặc vôn kế là cực kỳ quan trọng. Nếu chọn thang đo quá nhỏ so với giá trị cần đo, thiết bị có thể bị hỏng do quá tải. Ampe kế luôn được mắc nối tiếp để đo cường độ dòng điện, còn vôn kế luôn được mắc song song để đo hiệu điện thế. Sử dụng dây dẫn có tiết diện nhỏ có thể gây nóng chảy hoặc đứt dây, không liên quan trực tiếp đến an toàn khi sử dụng dụng cụ đo điện.
}
\end{ex}

% ===== Câu 45 =====
% ID: L10C1B2-124-TN
% Mức: VD
\begin{ex}
% ID: L10C1B2-124-TN
% Mức: VD
Khi sử dụng ampe kế để đo cường độ dòng điện trong mạch, thao tác nào sau đây là an toàn và đúng?
\choice
{Mắc ampe kế song song với nguồn điện}
{\True Chọn thang đo phù hợp và mắc ampe kế nối tiếp với mạch}
{Chạm tay trực tiếp vào hai đầu dây dẫn đang có điện}
{Tăng điện áp nguồn đến mức lớn nhất trước khi đo}
\loigiai{
Ampe kế phải mắc nối tiếp với mạch cần đo. Cần chọn thang đo phù hợp để tránh làm hỏng dụng cụ và bảo đảm an toàn.
}
\end{ex}

% ===== Câu 46 =====
% ID: L10C1B2-125-TN
% Mức: VD
\begin{ex}
% ID: L10C1B2-125-TN
% Mức: VD
Khi đun nóng chất lỏng trong ống nghiệm, cách làm nào sau đây là an toàn?
\choice
{Hướng miệng ống nghiệm về phía người quan sát}
{Đậy kín miệng ống nghiệm khi đang đun}
{\True Hướng miệng ống nghiệm ra xa người và hơ nóng đều}
{Cầm trực tiếp ống nghiệm bằng tay}
\loigiai{
Khi đun chất lỏng trong ống nghiệm, cần hướng miệng ống nghiệm ra xa người để tránh chất lỏng bắn ra gây bỏng.
}
\end{ex}

% ===== Câu 47 =====
% ID: L10C1B2-126-TN
% Mức: VD
\begin{ex}
% ID: L10C1B2-126-TN
% Mức: VD
Trong quá trình lắp đặt một mạch điện đơn giản để nghiên cứu về định luật Ôm, bạn Nam vô tình làm rơi một dụng cụ kim loại nhỏ (như tua vít) vào khu vực có các dây dẫn đang nối với nguồn điện có hiệu điện thế $12\,\text{V}$. Hiện tượng nào sau đây có khả năng xảy ra nhất và hành động nào là KHÔNG an toàn?
\choice
{Dụng cụ kim loại bị nóng lên nhanh chóng, có thể gây bỏng nhẹ cho Nam nếu chạm vào. Nam nên dùng tay nhấc dụng cụ ra khỏi mạch điện ngay lập tức.}
{\True Dụng cụ kim loại có thể gây ngắn mạch, làm tăng cường độ dòng điện đột ngột, có thể làm hỏng nguồn điện hoặc dây dẫn. Nam nên ngắt ngay nguồn điện trước khi xử lý.}
{Dụng cụ kim loại sẽ không bị ảnh hưởng gì vì hiệu điện thế chỉ 12 $V$ là quá nhỏ. Nam có thể thoải mái nhặt dụng cụ lên.}
{Dụng cụ kim loại sẽ bị chảy ra do nhiệt lượng sinh ra từ dòng điện. Nam nên dùng tay gạt dụng cụ ra khỏi mạch.}
\loigiai{
Khi một vật dẫn điện chạm vào hai điểm có chênh lệch điện thế trong mạch điện, nó tạo ra một đường dẫn cho dòng điện chạy qua. Nếu vật này nối trực tiếp hai cực của nguồn điện hoặc hai điểm có điện trở rất nhỏ, nó sẽ gây ra hiện tượng ngắn mạch. Ngắn mạch làm cho cường độ dòng điện tăng lên rất lớn, vượt quá khả năng chịu đựng của dây dẫn và nguồn điện, dẫn đến sinh nhiệt lớn, có thể làm cháy dây, hỏng thiết bị hoặc thậm chí gây hỏa hoạn. Hành động an toàn nhất trong trường hợp này là ngắt nguồn điện ngay lập tức để loại bỏ nguy cơ gây hại trước khi xử lý vật thể lạ trong mạch.
}
\end{ex}

% ===== Câu 124 =====
% ID: L10C1B2-39-DS
% Mức: NB
\begin{ex}
% ID: L10C1B2-39-DS
% Mức: NB
Những hoạt động tuân thủ nguyên tắc an toàn khi sử dụng điện:
\choiceTF
{\True Bọc kĩ các dây dẫn điện bằng vật liệu cách điện}
{\True Kiểm tra mạch có điện bằng bút thử điện}
{Sửa chữa điện khi chưa ngắt nguồn điện}
{Chạm tay trực tiếp vào ổ điện, dây điện trần hoặc dây dẫn điện bị hở}
\loigiai{
\begin{itemchoice}
\item (Đúng) Bọc kĩ các dây dẫn điện bằng vật liệu cách điện. (Vì): Việc bọc dây dẫn điện bằng vật liệu cách điện là biện pháp cơ bản để ngăn cách người sử dụng với dòng điện, từ đó phòng tránh nguy cơ bị điện giật
\item (Đúng) Kiểm tra mạch có điện bằng bút thử điện. (Vì): Sử dụng bút thử điện để kiểm tra xem mạch điện có còn điện hay không trước khi tiến hành sửa chữa là một bước quan trọng để đảm bảo an toàn, tránh tai nạn đáng tiếc
\item (Sai) Sửa chữa điện khi chưa ngắt nguồn điện. (Vì): ửa chữa điện khi nguồn điện vẫn đang hoạt động là hành động cực kỳ nguy hiểm, tiềm ẩn nguy cơ cao bị điện giật, đe dọa trực tiếp đến tính mạng
\item (Sai) Chạm tay trực tiếp vào ổ điện, dây điện trần hoặc dây dẫn điện bị hở. (Vì): Tiếp xúc trực tiếp với các bộ phận mang điện hở như ổ điện, dây dẫn trần hoặc dây bị hở vỏ cách điện sẽ gây ra hiện tượng điện giật, rất nguy hiểm cho sức khỏe và tính mạng
\end{itemchoice}
}
\end{ex}

% ===== Câu 127 =====
% ID: L10C1B2-40-DS
% Mức: NB
\begin{ex}
% ID: L10C1B2-40-DS
% Mức: NB
Trong phòng thí nghiệm, cần lưu ý:
\choiceTF
{\True Không ăn uống trong phòng thí nghiệm}
{\True Để hóa chất chung với thực phẩm để tiết kiệm không gian}
{\True Đeo kính bảo hộ khi làm việc với hóa chất}
{Rửa tay bằng xăng để loại bỏ hóa chất}
\loigiai{
\begin{itemchoice}
\item (Đúng) Không ăn uống trong phòng thí nghiệm. (Vì): Không nên ăn uống trong phòng thí nghiệm để tránh nguy cơ nhiễm độc hoặc tiếp xúc với hóa chất
\item (Đúng) Để hóa chất chung với thực phẩm để tiết kiệm không gian. (Vì): Không được để hóa chất chung với thực phẩm vì dễ gây nguy hiểm cho sức khỏe
\item (Đúng) Đeo kính bảo hộ khi làm việc với hóa chất. (Vì): eo kính bảo hộ giúp bảo vệ mắt khi làm việc với hóa chất hoặc các chất nguy hiểm
\item (Sai) Rửa tay bằng xăng để loại bỏ hóa chất. (Vì): Rửa tay bằng xăng có thể gây tổn thương da và không loại bỏ hoàn toàn hóa chất, nên sử dụng xà phòng và nước để rửa tay
\end{itemchoice}
}
\end{ex}

% ===== Câu 129 =====
% ID: L10C1B2-41-DS
% Mức: NB
\begin{ex}
% ID: L10C1B2-41-DS
% Mức: NB
Khi sử dụng máy móc trong phòng thí nghiệm:
\choiceTF
{\True Đọc kỹ hướng dẫn sử dụng trước khi dùng}
{Dùng máy ngay lập tức mà không cần kiểm tra}
{\True Để máy gần nguồn nước để tiện sử dụng}
{\True Tắt máy khi không sử dụng}
\loigiai{
\begin{itemchoice}
\item (Đúng) Đọc kỹ hướng dẫn sử dụng trước khi dùng. (Vì): ọc kỹ hướng dẫn sử dụng giúp hiểu rõ cách thức vận hành, các chức năng, giới hạn hoạt động và các biện pháp phòng ngừa, từ đó tránh các lỗi hoặc tai nạn không đáng có trong quá trình sử dụng
\item (Sai) Dùng máy ngay lập tức mà không cần kiểm tra. (Vì): Việc không kiểm tra máy trước khi sử dụng có thể dẫn đến các sự cố kỹ thuật, hỏng hóc thiết bị hoặc gây nguy hiểm cho người sử dụng và môi trường xung quanh. Cần kiểm tra sơ bộ tình trạng máy, các kết nối và các bộ phận an toàn trước khi vận hành
\item (Đúng) Để máy gần nguồn nước để tiện sử dụng. (Vì): ể máy móc, đặc biệt là các thiết bị điện, gần nguồn nước có thể gây nguy hiểm do chập điện, giật điện, làm hỏng máy và tiềm ẩn nguy cơ cháy nổ
\item (Đúng) Tắt máy khi không sử dụng. (Vì): Tắt máy khi không sử dụng không chỉ giúp tiết kiệm năng lượng điện tiêu thụ mà còn đảm bảo an toàn, tránh các sự cố có thể xảy ra khi máy hoạt động không cần thiết hoặc khi có sự cố về điện
\end{itemchoice}
}
\end{ex}

% ===== Câu 135 =====
% ID: L10C1B2-44-DS
% Mức: NB
\begin{ex}
% ID: L10C1B2-44-DS
% Mức: NB
Khi sử dụng dao cắt trong phòng thí nghiệm:
\choiceTF
{\True Đeo găng tay bảo hộ để tránh bị cắt}
{\True Cắt trực tiếp về phía mình để dễ điều khiển}
{\True Sử dụng dao cắt sắc để dễ dàng cắt chính xác}
{Dùng tay không để giữ vật cần cắt}
\loigiai{
\begin{itemchoice}
\item (Đúng) Đeo găng tay bảo hộ để tránh bị cắt. (Vì): eo găng tay bảo hộ giúp bảo vệ tay khỏi bị cắt và đảm bảo an toàn hơn khi thao tác
\item (Đúng) Cắt trực tiếp về phía mình để dễ điều khiển. (Vì): Cắt trực tiếp về phía mình là hành động rất nguy hiểm, có thể gây tổn thương nghiêm trọng cho người sử dụng
\item (Đúng) Sử dụng dao cắt sắc để dễ dàng cắt chính xác. (Vì): Sử dụng dao cắt sắc giúp thao tác chính xác, giảm lực tác động và tránh tai nạn do lưỡi dao cùn gây ra
\item (Sai) Dùng tay không để giữ vật cần cắt. (Vì): Dùng tay không để giữ vật cắt là hành động cực kỳ nguy hiểm. Cần sử dụng kẹp hoặc các dụng cụ hỗ trợ chuyên dụng để giữ vật cần cắt
\end{itemchoice}
}
\end{ex}

% ===== Câu 136 =====
% ID: L10C1B2-45-DS
% Mức: TH
\begin{ex}
% ID: L10C1B2-45-DS
% Mức: TH
Khi tiến hành phản ứng hóa học, cần tuân thủ quy tắc:
\choiceTF
{Đổ nước vào axit để làm loãng axit}
{\True Đổ axit từ từ vào nước để tránh bắn tung tóe}
{Đổ axit vào bột kim loại để tăng tốc độ phản ứng}
{\True Làm việc trong tủ hút để tránh hít phải khí độc}
\loigiai{
\begin{itemchoice}
\item (Sai) Đổ nước vào axit để làm loãng axit. (Vì): Khi pha loãng axit, quy tắc an toàn là phải đổ từ từ axit đặc vào nước, không bao giờ đổ ngược lại. Việc đổ nước vào axit đặc có thể gây ra phản ứng tỏa nhiệt mạnh, làm sôi dung dịch và bắn axit ra ngoài, gây bỏng nặng
\item (Đúng) Đổ axit từ từ vào nước để tránh bắn tung tóe. (Vì): ổ từ từ axit đặc vào nước giúp nhiệt lượng tỏa ra trong quá trình pha loãng được nước hấp thụ và phân tán dần, tránh làm dung dịch sôi đột ngột và bắn tung tóe. Điều này giúp kiểm soát phản ứng và đảm bảo an toàn
\item (Sai) Đổ axit vào bột kim loại để tăng tốc độ phản ứng. (Vì): Đổ axit vào bột kim loại có thể gây ra phản ứng hóa học mạnh, tỏa nhiều nhiệt và sinh ra khí độc, tùy thuộc vào loại axit và kim loại. Việc này tiềm ẩn nguy cơ cháy nổ hoặc ngộ độc nếu không được thực hiện trong điều kiện kiểm soát chặt chẽ và có biện pháp phòng ngừa phù hợp
\item (Đúng) Làm việc trong tủ hút để tránh hít phải khí độc. (Vì): Làm việc trong tủ hút là một biện pháp an toàn quan trọng khi thực hiện các phản ứng hóa học có khả năng sinh ra khí độc, hơi độc hoặc bụi độc hại. Tủ hút giúp ngăn chặn sự phát tán của các chất nguy hiểm ra môi trường làm việc, bảo vệ sức khỏe và an toàn cho người thực hiện thí nghiệm
\end{itemchoice}
}
\end{ex}

% ===== Câu 138 =====
% ID: L10C1B2-46-DS
% Mức: TH
\begin{ex}
% ID: L10C1B2-46-DS
% Mức: TH
Các qui tắc an toàn khi sử dụng thiết bị thí nghiệm:
\choiceTF
{Sử dụng thiết bị điện: Sử dụng điện áp lớn hơn thông số qui định trên thiết bị để thiết bị hoạt động mạnh hơn}
{\True Sử dụng thiết bị nhiệt và thủy tinh: Các thiết bị nung nóng có thể gây cháy hoặc nứt vỡ các bộ phận làm thủy tinh}
{Sử dụng các thiết bị quang học: Các thiết bị quang học rất dễ mốc, xước, nứt, vỡ và dính bụi bẩn, làm ảnh hưởng đến đường truyền tia sáng và sai lệch kết quả thí nghiệm}
{\True Tuân thủ quy tắc an toàn phòng cháy chữa cháy và an toàn khi sử dụng hóa chất dễ cháy, nổ}
\loigiai{
\begin{itemchoice}
\item (Sai) Sử dụng thiết bị điện: Sử dụng điện áp lớn hơn thông số qui định trên thiết bị để thiết bị hoạt động mạnh hơn. (Vì): Khi pha loãng axit, phải đổ axit vào nước từ từ để tránh phản ứng mạnh, không đổ nước vào axit
\item (Đúng) Sử dụng thiết bị nhiệt và thủy tinh: Các thiết bị nung nóng có thể gây cháy hoặc nứt vỡ các bộ phận làm thủy tinh. (Vì): ổ axit vào nước từ từ giúp kiểm soát phản ứng và tránh tình trạng bắn tung tóe, gây nguy hiểm
\item (Sai) Sử dụng các thiết bị quang học: Các thiết bị quang học rất dễ mốc, xước, nứt, vỡ và dính bụi bẩn, làm ảnh hưởng đến đường truyền tia sáng và sai lệch kết quả thí nghiệm. (Vì): Đổ axit vào bột kim loại có thể gây phản ứng mạnh, nguy hiểm. Không nên làm điều này mà không có hướng dẫn rõ ràng
\item (Đúng) Tuân thủ quy tắc an toàn phòng cháy chữa cháy và an toàn khi sử dụng hóa chất dễ cháy, nổ. (Vì): Làm việc trong tủ hút giúp tránh hít phải khí độc, bảo vệ sức khỏe và an toàn cho người thực hiện
\end{itemchoice}
}
\end{ex}

% ===== Câu 140 =====
% ID: L10C1B2-47-DS
% Mức: TH
\begin{ex}
% ID: L10C1B2-47-DS
% Mức: TH
An toàn trong phòng thực hành vật lí đòi hỏi sự cẩn trọng và hiểu biết về cách ứng phó với các tình huống nguy hiểm. Hãy đánh giá tính đúng/sai của các phát biểu sau đây.
\choiceTF
{Khi đo cường độ dòng điện hoặc hiệu điện thế, việc chọn thang đo nhỏ hơn giá trị dự kiến sẽ giúp tăng độ chính xác của phép đo.}
{\True Nếu phát hiện có mùi lạ hoặc khói bốc ra từ thiết bị điện đang hoạt động, cần ngắt nguồn điện ngay lập tức và báo cáo giáo viên.}
{Để tránh làm vỡ nhiệt kế thủy ngân, nên dùng tay không cầm trực tiếp vào bầu nhiệt kế khi đang đo nhiệt độ cao.}
{\True Khi làm việc với các thí nghiệm có vật nặng treo hoặc chuyển động, cần đảm bảo khu vực xung quanh không có người đứng gần để tránh va đập.}
\loigiai{
\begin{itemchoice}
\item (Sai) Khi đo cường độ dòng điện hoặc hiệu điện thế, việc chọn thang đo nhỏ hơn giá trị dự kiến sẽ giúp tăng độ chính xác của phép đo. (Vì): Chọn thang đo nhỏ hơn giá trị dự kiến có thể làm hỏng thiết bị đo do quá tải hoặc gây nguy hiểm. Cần chọn thang đo lớn hơn hoặc bằng giá trị dự kiến, sau đó điều chỉnh xuống thang đo phù hợp nhất để đạt độ chính xác cao
\item (Đúng) Nếu phát hiện có mùi lạ hoặc khói bốc ra từ thiết bị điện đang hoạt động, cần ngắt nguồn điện ngay lập tức và báo cáo giáo viên. (Vì): Mùi lạ hoặc khói là dấu hiệu của sự cố điện
\item (Sai) Để tránh làm vỡ nhiệt kế thủy ngân, nên dùng tay không cầm trực tiếp vào bầu nhiệt kế khi đang đo nhiệt độ cao. (Vì): Dùng tay không cầm trực tiếp vào bầu nhiệt kế khi đo nhiệt độ cao có thể gây bỏng cho người thực hiện. Ngoài ra, nhiệt độ từ tay có thể ảnh hưởng đến kết quả đo. Để tránh vỡ nhiệt kế, cần cẩn thận khi thao tác và tránh va đập mạnh
\item (Đúng) Khi làm việc với các thí nghiệm có vật nặng treo hoặc chuyển động, cần đảm bảo khu vực xung quanh không có người đứng gần để tránh va đập.
\end{itemchoice}
}
\end{ex}

% ===== Câu 142 =====
% ID: L10C1B2-48-DS
% Mức: TH
\begin{ex}
% ID: L10C1B2-48-DS
% Mức: TH
Việc sử dụng đúng cách các thiết bị bảo hộ và hiểu rõ các nguy cơ tiềm ẩn là chìa khóa để đảm bảo an toàn trong phòng thí nghiệm vật lí. Hãy đánh giá tính đúng/sai của các phát biểu sau.
\choiceTF
{\True Luôn đeo kính bảo hộ khi thực hiện các thí nghiệm có nguy cơ bắn vật nhỏ, hóa chất hoặc ánh sáng mạnh để bảo vệ mắt.}
{\True Không được tự ý sửa chữa hoặc tháo rời các thiết bị điện khi chúng đang được cắm vào nguồn điện, ngay cả khi chúng không hoạt động.}
{Trong trường hợp bị điện giật nhẹ, chỉ cần rút phích cắm ra khỏi ổ điện là đủ, không cần thiết phải báo cáo giáo viên.}
{Khi cần di chuyển các thiết bị nặng hoặc cồng kềnh, nên tự mình thực hiện để không làm phiền người khác.}
\loigiai{
\begin{itemchoice}
\item (Đúng) Luôn đeo kính bảo hộ khi thực hiện các thí nghiệm có nguy cơ bắn vật nhỏ, hóa chất hoặc ánh sáng mạnh để bảo vệ mắt. (Vì): Mắt là bộ phận nhạy cảm và dễ bị tổn thương. Kính bảo hộ là thiết bị thiết yếu để bảo vệ mắt khỏi các tác nhân gây hại như mảnh vỡ, hóa chất bắn vào hoặc ánh sáng cường độ cao
\item (Đúng) Không được tự ý sửa chữa hoặc tháo rời các thiết bị điện khi chúng đang được cắm vào nguồn điện, ngay cả khi chúng không hoạt động. (Vì): Ngay cả khi thiết bị không hoạt động, nếu vẫn được cắm vào nguồn điện, vẫn có nguy cơ bị điện giật do lỗi bên trong hoặc chạm vào các bộ phận mang điện. Mọi hoạt động sửa chữa, tháo lắp phải được thực hiện khi thiết bị đã ngắt hoàn toàn khỏi nguồn điện và dưới sự giám sát của giáo viên
\item (Sai) Trong trường hợp bị điện giật nhẹ, chỉ cần rút phích cắm ra khỏi ổ điện là đủ, không cần thiết phải báo cáo giáo viên. (Vì): Mọi trường hợp bị điện giật, dù nhẹ, đều cần được báo cáo ngay lập tức cho giáo viên và cần được kiểm tra y tế. Điện giật có thể gây ra những tổn thương nội tạng không thể nhìn thấy ngay lập tức
\item (Sai) Khi cần di chuyển các thiết bị nặng hoặc cồng kềnh, nên tự mình thực hiện để không làm phiền người khác. (Vì): Di chuyển thiết bị nặng hoặc cồng kềnh một mình có thể gây chấn thương cho bản thân (như đau lưng, trượt chân) hoặc làm hỏng thiết bị. Cần có sự hỗ trợ của người khác hoặc sử dụng các dụng cụ hỗ trợ di chuyển, và phải có sự hướng dẫn của giáo viên
\end{itemchoice}
}
\end{ex}

